\documentclass[a4paper,12pt]{article}

\usepackage[utf8]{inputenc}
\usepackage[T1]{fontenc}
\usepackage[swedish]{babel}
\usepackage{geometry}
\geometry{margin=2.5cm}
\usepackage{setspace}
\onehalfspacing
\usepackage{parskip}
\setlength{\parindent}{0pt}
\usepackage{hyperref}
\usepackage{graphicx}
\usepackage{booktabs}
\usepackage{array}

\begin{document}

% ================== FRAMSIDA ==================
\begin{titlepage}
    \centering

    {\Large TDDE02 – Mjukvarutekniskt entreprenörskap\par}
    \vspace{2cm}

    {\LARGE \textbf{Projektrapport – Affärsanalys}}\par
    \vspace{1.5cm}

    {\large Projekt: \par}
    \vspace{1cm}

    {\large Grupp 5\par}
    \vspace{0.5cm}

   
    {\small

    Jonatan Ebenholm (joneb274)\\
   Leon Sandenskog (leosa068)\\
   Bushra Alawad (busal420)\\
   Berkay Orhan (beror658)\\
   Hani Abdallah (hanab049)
    }\par

    \vfill
    {\large \par}

\end{titlepage}


\newpage
\tableofcontents

\newpage

% ================== 1. SAMMANFATTNING ==================

\section{Sammanfattning}
% Kort: vad idén är, vad analysen visar, rekommendation (GO / KILL / RECYCLE).


% ================== 2. BAKGRUND OCH AFFÄRSIDÉ ==================
\section{Bakgrund och översiktlig beskrivning av affärsidén}
% Problem/utmaning, data som visar att behovet finns, avsluta med vision.
% Beskriv affärsidén enligt NABC (ev. flera NABC om flera kundtyper).

\subsection{Bakgrund och problem}
Varje dag står många av oss i köket och tittar in i ett fullt kylskåp, men ändå har vi ingen aning om vad vi ska laga till middag. Det är ett vardagsproblem som leder till onödig stress, onödigt matsvinn och i slutändan slösade pengar.
Siffrorna talar sitt tydliga språk att varje person i Sverige slänger i genomsnitt 60 kilo mat per år(källa). Ungefär 16 kilo av det som kastas är faktiskt fortfarande helt ätbart. Samtidigt lägger svenska hushåll mer än 300 miljarder kronor om året på maten.

Många försöker hitta sätt att spara tid, minska slöseri och få mer kontroll över matlagen. Det är precis där vår tjänst kommer in. Vi vill göra det enkelt att använda det man redan har hemma så att planeringen känns lättare, maten räcker längre och mindre slängs i soporna.


\subsection{Behov och data}
Vi gjorde en kundundersökning som omfattade 23 personer i åldern 21–30 år. Utifrån kundundersökningen(källa) har 52\% uppger att de ofta eller alltid har svårt att komma på vad de ska laga, 69,6 \% vill ha recept baserade på ingredienser de redan har hemma, 73,9 \% vill kunna använda kameran för att registrera sina ingredienser och endast 13 \% säger att de aldrig slänger mat.

Både nationell statistik och vår egen undersökning visar att behovet är stort. Det finns en tydlig efterfrågan på ett verktyg som kan ge struktur, minska matsvinn och förenkla vardagen utan att användaren behöver lägga extra tid på att lista ingredienser eller söka recept manuellt.


\subsection{Vision}
Vi har en enkel vision. Det är att vi vill göra matplaneringen enkelt och smidig med hjälp av smart teknik. Tanken är att du ska kunna använda det du redan har hemma, istället för att det ska bli över och slängas. På så sätt slipper du slösa mat, tid och pengar.
I praktiken innebär det att du ska kunna öppna kylskåpet, ta några bilder och direkt få tips på maträtter som du faktiskt har ingredienser till. Ingen stress, ingen onödig shopping, ingen mat som blir kvar och glöms bort.
Långsiktigt drömmer vi om att bli den naturliga digitala köksassistenten i svenska hem och med tiden kanske även utanför Sveriges gränser.



\subsection{Affärsidén enligt NABC}
% Kort NABC-översikt, detaljer kan komma senare i Erbjudandet.

-Need (Behov): Undersökningen och i den nationella statistiken -visar att många kämpar med planeringen kring mat. Det är inte alltid bristen på mat, utan mer brist på idéer. Många vill undvika att slänga mat och därför behövs det hjälp i vardagen för att faktiskt lyckas. Dessutom vill en del av användarna ha stöd i själva matlagningen, till exempel genom att få rekommenderade instruktionsvideos.
-Approach (Lösning) :Appen ska göra det lätt att få koll på vad man har hemma och hur man kan använda det. Appen bygger på tre grundidéer:

1- Ta bilder på kylskåpet från flera vinklar för att få med alla ingredienser utan att behöva skriva något själv.

2- Automatiska receptförslag som baseras på ingredienserna och som går att anpassa utifrån tid, kostval eller hur avancerat man vill ha det.

3- Fokus på minimal komplettering, alltså recept där man redan har nästan allt.

-Benefit (Fördelar): För användaren handlar fördelarna om att få en enklare vardag. Mindre tid på att tänka ut vad man ska laga, färre impulsköp och mindre mat som hinner bli gammal. På ett större plan kan lösningen bidra till att minska matsvinnet och därmed också spara pengar och miljöpåverkan.
För företag och samarbetspartners finns möjligheter till premiumfunktioner, samarbeten och ytterligare affärsmodeller.

-Competition (Konkurrens): Visst finns det massor med receptappar, men många kräver att du själv skriver in allt du har hemma. Andra är mest till för inspiration. De visar fina bilder och videor, men hjälper dig inte när du faktiskt ska laga mat.
Till exempel kräver Supercook att du matar in varje ingrediens för hand. Tasty är kul att titta på, men ger inga förslag baserat på ditt eget kylskåp. Recept.se är som ett enormt digitalt receptarkiv. Du får själv lista ut vad du ska söka på.
Vi erbjuder dig att kunna använda det som du redan har i ditt kylskåp. Istället för att bara visa massor recept, hjälper vi dig att få konkreta förslag som passar just din vardag, dina ingredienser och din tid.





% ================== 3. ERBJUDANDET ==================
\section{Erbjudandet}
% Mer detaljerad beskrivning av lösning, teknik, skalbarhet, IPR, Lean Canvas-översikt.

\subsection{Värdeerbjudande}
De flesta har varit i situationen där man står framför ett kylskåp fullt med mat men ändå inte vet vad man ska laga. Ofta handlar det inte om att man saknar ingredienser, utan att det känns svårt att få ihop något av det man har. Det är just det problemet som ska lösas. Genom appen ska man enkelt kunna ta några bilder av kylskåpet. Appen identifierar ingredienserna och föreslår rätter som går att laga utan massa extra inköp. Tjänsten riktar sig främst till unga vuxna och andra som vill laga mer mat hemma men som saknar tid eller inspiration. Tanken är att appen inte ska vara ännu en receptsamling, utan ett vardagsverktyg som hjälper till att hålla koll, minska stress och få ut mer av maten man redan har.


\subsection{Produkt/tjänst i moget läge}
När appen är fullt utvecklad ska den fungera som en hjälpreda i köket. Den ska kunna hantera flera bilder och sätta ihop en ingredienslista. Den ska känna igen vanliga livsmedel automatiskt. Plus att den kan föreslå några recept som matchar det användaren har. Dessutom låter den användaren filtrera recept efter tid, nivå, diet eller allergier. Den kan även tipsa om instruktionsvideos för den som vill ha extra guidning. Om ingredienserna saknas kommer appen att föreslå rätter som kräver så få kompletteringar som möjligt. Användare kan också spara sina favoritrecept i applikationen.


\subsection{Teknik, skalbarhet och genomförbarhet}

Den tekniska lösningen bygger på att vi använder bildigenkänning från befintliga API:er och en mobilapp som fungerar på både Android och iPhone. I appen finns även en AI modell som kan sätta ihop recept på ett tydligt och förståeligt sätt för svenska användare
Skalbarheten är god eftersom det främst är serverkapacitet som kostar när antalet användare växer, inte genom utvecklingstiden. Dessutom blir tekniken billigare över tid.
Med hjälp av Immaterialrättsligt kan vi skydda:

-algoritmen som matchar ingredienser: Det betyder att ingen annan kan använda exakt samma algoritm utan tillstånd.


-multibild: sättet appen hanterar flera bilder för att skapa en komplett ingredienslista kan också vara skyddat.


-designen och användarupplevelsen: hur appen eller menyer ser ut  och hur användaren interagerar med den kan skyddas genom upphovsrätt eller designskydd.


-Varumärket: såsom namnet på appen, logotypen och andra symboler som identifierar tjänsten.

\subsection{Affärsmodell och Lean Canvas – översikt}

Vår affärsmodell bygger på att appen är gratis i grunden, men att ha extra funktioner kräver ett premiumabonnemang. Vi vill att vår målgrupp ska vara unga vuxna, personer med lite tid, och de som vill minska matsvinn eller få struktur som underlättar vardagen. Vår nyckelkompetens handlar om teknik, användarupplevelse, receptlogik och livsmedelskunskap. Våra viktiga partners kan vara matkedjor, API-leverantörer, matinfluenser och livsmedelsproducenter. För att kunna försörja vår process behöver vi ha intäkter. Våra intäkterna kan komma från, premium, samarbeten, sponsrat innehåll och anonymiserad marknadsdata. Utifrån Lean Canvas visar det att vår lösning har en tydlig värdeproposition, realistiska kostnader, flera möjliga intäktskällor och en målgrupp som både behöver och aktivt efterfrågar tjänsten.


% ================== 4. OMVÄRLDSANALYS ==================
\section{Omvärldsanalys}
% Bransch, konkurrenter, andra aktörer. Styrkor/svagheter hos konkurrenter.

\subsection{Branschbeskrivning}
Vår app hamnar egentligen mitt i mitten mellan mat, teknik och hållbarhet. Vi rör oss inom områden som digitala recept, matlagningsappar och smarta kök assistenter.

Utvecklingen inom det här området har gått väldigt snabbt de senaste åren. I dag planerar många sin mat digitalt, handlar via mobilen och använder appar för både recept och inspiration. Samtidigt är hållbarhet något som blivit en självklar del av vardagen för många.

Trots det finns det fortfarande få tjänster som på ett riktigt smart sätt använder mobilkameran för att hjälpa till med själva matplaneringen.Du får massa förslag på vad du skulle kunna laga, men ingen som visar hur du ska använda just det du har hemma.

Det är just det här tomrummet vi vill fylla. Med vår lösning vill vi göra vardagen lite enklare, hjälpa människor att ta bättre vara på sin mat och samtidigt bidra till att minska onödigt matsvinn.

\subsection{Konkurrenter och substitut}
Det finns många appar och webbplatser för recept, men de fungerar ganska olika jämfört med det vi vill erbjuda. Här är några av de mest relevanta konkurrenterna och hur de skiljer sig från oss.

Supercook:

Supercook är ett stort namn som många känner till. Där skriver du in allt du har hemma och får receptförslag. Problemet är att det är du själv som måste skriva in allting och det tar tid. Efter en lång dag orkar man sällan sitta och lista upp varje sak i kylskåpet. De har ingen bildigenkänning och recepten är inte anpassade efter hur mycket tid du har eller vad du föredrar att äta.

Det Supercook gör bra är att de har otroligt många recept. Men det känns ändå lite platt för allt börjar med att du själv ska fylla i en lista.

Vår idé är att det ska räcka med att du tar några bilder av kylskåpet. Inget manuellt skrivande, bara enkelt och smidigt. Vi vill att förslagen ska passa just för din situation. Om du har bråttom, om du saknar en ingrediens eller om du vill äta på ett visst sätt.

Tasty (BuzzFeed):

Tasty är super populärt. De har snygga videor, roliga idéer och man kan fastna länge och bara titta. Men de är mer underhållning än verktyg. Det finns ingen koppling till vad du faktiskt har hemma. Du får inspiration, men inte hjälp med att planera din egen matlagning.

Vi tänkte göra något som faktiskt hjälper på riktigt. Vår app utgår från just ditt kylskåp. Alltså ,vad du redan har och den ger dig matförslag du faktiskt kan laga nu på en gång.

Recept.se:

Recept.se är som en jättestor digital kokbok med massor av recept. Men du måste veta vad du ska söka på och det finns inget som hjälper dig hålla kolla på det som kommer bli gammalt. Du får ingen hjälp med att undvika att köpa saker i onödan, eller att faktiskt använda det du redan har hemma.

Vi vill vara som din personliga köksassistent. Appen ska utgå från ditt kylskåp istället för en stor lista med recept. Tanken är att den ska hjälpa dig hela vägen. Från att se vad du har till att välja recept och undvika onödiga inköp.

\subsection{Sammanfattande slutsatser från omvärldsanalysen}
Utifrån omvärldsanalysen har vi kommit fram till att behovet redan finns och målgruppen är ganska stor. Vår lösning är unikt jämför med de andra aktörerna. De erbjuder bara vissa delar av det vi vill göra och har. Marknaden är inte full ännu och därför har vi möjligheten att konkurrera i den. Vår största risk är att ett större företag skulle kunna bygga något liknande snabbt om de ville.

Allt detta ger bra underlag inför både SWOT-analysen och Lean Canvas. Det stärker också argumenten för varför vår lösning har en tydlig konkurrensfördel och faktiskt har chans att lyckas.


% ================== 5. MARKNADSANALYS ==================
\section{Marknadsanalys}
% Val av målgrupp, early adopters, kanaler, kundrelationer, empiriska studier, MVP-tester.

\subsection{Målgrupp och early adopters}


\subsection{Marknadsstorlek och antaganden}


\subsection{Kanaler och kundrelationer}


\subsection{Empiriska studier (intervjuer, enkät, MVP-test)}


\subsection{Sammanfattande slutsatser från marknadsanalysen}


% ================== 6. ORGANISATION OCH TEAM ==================
\section{Organisation och team}
% Framtida organisation, värdekedja, aktiviteter, resurser, partners, teamroller.

\subsection{Framtida organisation och värdekedja}


\subsection{Nyckelaktiviteter och resurser}


\subsection{Team och kompetenser}


\subsection{Partners och externa resurser}


% ================== 7. ANALYS AV BÄRKRAFT ==================
\section{Analys av bärkraften samt lönsamhet och finansiering}

\subsection{SWOT-analys}
% Styrkor, svagheter, möjligheter, hot.


\subsection{Affärsmodell}
% Beskriv vald affärsmodell, koppla till analyserna.


\subsection{Ekonomi och finansiering}
% Rimliga antaganden om intäkter, kostnader, betalvilja, finansieringsvägar.


% ================== 8. GROWTH HACKING ==================
\section{Growth hacking}
% Funktioner och strategier för spridning, viralitet, inbyggd marknadsföring.
För att snabbt etablera appen på marknaden utan kostsamma traditionella reklamkampanjer, tillämpas en strategi baserad på growth hacking. Fokus ligger på att utnyttja produktens \textit{Unika värdeerbjudande}; enkelheten i att ta en bild och få ett recept, vilket kan hjälpa att skapa organisk spridning.

Eftersom våra primära kundsegment består av pris- och tidsmedvetna grupper såsom studenter och gym-bros, är målet att integrera marknadsföringen direkt i användarupplevelsen. Genom att fokusera på våra identifierade kanaler, främst sociala medier och \textit{word-of-mouth}, ska vi konvertera användare till ambassadörer. Nedan presenteras strategier för att driva snabb tillväxt genom lågkostnadsinitiativ och community-byggande.

\subsection{Challanges}
Eftersom appen bygger på att ta bilder på ingredienser så skulle en TikTok/Instagram kampanj skapas då användare som exempel kan ta en bild på sitt kylskåp och sedan sin måltid som appen producerat recept för. Det skulle vara lite som en \textit{before and after} jämförelse. Det är ett visuellt och delbart sätt att engagera personer och få gratis reklam via sociala medier.

\subsection{Ambassadörsprogram för Studenter och Gymbros} 
Då vi tror att studenter och gymbros är bland de tydligaste \textit{early adopters} vi kommer ha så kommer vi aktivt söka deras engagemang genom att göra det möjligt för dem att agera som ambassadörer och rekrytera personer till appen. Detta skulle fungera så att om en person rekryterar säg fem personer till appen så får de förmåner så som premiumfunktioner. Där efter så kanske de får premiumfunktionerna på längre sikt om de forsätter rekrytera. Alla som har ett konto är ambassadörer, men det finns inget krav att rekrytera, endast de som söker fler funktioner utan att betala behöver rekrytera. Om användarna de rekryterar betalar för premium eller rekryterar vidare ännu fler personer så kan detta också gynna den tidigare nämnda ambassadören. Detta kommer kunna leda till att användare vi är rätt säkra på kommer gilla appen kommer att också vilja dela med sig av den till sina vänner och kontakter.

\subsection{Partnerskap med studentkårer}
Partnerskap mellan studentkårer under till exempel nolle-p kan locka studenter som kanske aldrig förr behövt tänka mycket på om saker och ting kommer gå ut i sitt kylskåp eller aldrig haft ont om pengar för mat. 

\subsection{Gamification}
Om vi mäter matsvinn så kan vi också mäta hur mycket matsvinn som har minskats, vilket kan användas för att ge användaren en mätare för både matsvinn reducerat och även pengar sparat. Om användaren sedan kan dela denna statistik så ger det folk en idé om att de kan betala för appen för att spara pengar i slutändan medan de också slipper åka en extra vända till mataffären. 




% ================== 9. FRAMTIDA UTVECKLINGSMÖJLIGHETER ==================
\section{Framtida utvecklingsmöjligheter}
% Nya kundsegment, teknisk utveckling, partners, expansion bortom första nischen.
Lanseringen av appen är endast första steget mot att skapa ett heltäckande ekosystem för smart matlagning. Medan den initiala versionen fokuserar på kärnproblemet, att matcha befintliga ingredienser med recept, så kräver långsiktig bärkraft en expansion av både funktionalitet och affärsmodell.

För att säkerställa att vi inte bara attraherar utan även behåller våra användare över tid, måste vi utveckla appen från ett enkelt verktyg till en oumbärlig del av hushållsekonomin och hälsan. Detta inkluderar att realisera våra identifierade intäktsströmmar, såsom betalversioner och B2B-samarbeten (Business-to-Business). I detta avsnitt redogörs för visionen framåt, med fokus på teknisk innovation och strategiska partnerskap som ökar appens värde för både konsumenter och samarbetspartners.

\subsection{Integration med E-handel}
Om vi i framtiden kan integrera B2B-samarbeten i appen så skulle det vara en inkomstström extra. Precis hur denna integration skulle se ut är inte klar, men det kan vara ett samarbete där om användaren skulle vilja så kan de koppla sin pentry till Ica, Coop, Lidl eller andra matkedjor för att generera recept med mat som är nedsatt hos matkedjorna. Detta skulle igen bidra till idén om att appen är ett sätt att spara pengar på.

\subsection{Community-driven AI-träning}
Tanken är att användaren ska kunna ändra på sina AI-genererade recept. Dessa ändringar skulle kunna lagras som en datapunkt där de kan användas för att träna AI-modellen till att ta bort krydda som de flesta användarna inte tycker om. Detta skulle också vara intressant data för att se om användare överlag ogillar specifika råvaror som restauranger eller butiker kanske hade varit intresserade av att känna till.


% ================== 10. HANDLINGS- / GENOMFÖRANDEPLAN ==================


\section{Genomförda iterationer och MVP:er}




% ================== 11. LEAN CANVAS ==================
\section{Lean Canvas}
% Här kan ni antingen beskriva varje ruta eller lägga in en figur.

\subsection{Beskrivning av Lean Canvas}


% Exempel om ni vill lägga in en bild:
% \begin{figure}[h]
%     \centering
%     \includegraphics[width=0.9\textwidth]{bilder/lean_canvas.png}
%     \caption{Lean Canvas för [projektnamn]}
% \end{figure}

% ================== 12. SLUTSATS ==================
\section{Slutsats}
% Sammanfattande analys och tydlig rekommendation: GO / KILL / RECYCLE.



% ================== REFERENSER ==================
\section*{Referenser}
% Om ni använder BibTeX, ersätt detta med \bibliographystyle + \bibliography.
% Annars skriv manuellt här.



% ================== BILAGOR ==================
\appendix






\end{document}
